%%% The 65/41 Barrier: Continued-Fraction Convergents of log_2 3
%%% Bound Worst-Case Descent Ratios in the Compressed Collatz Map

\documentclass[11pt,a4paper]{article}

%% ---------- packages ----------
\usepackage[utf8]{inputenc}
\usepackage{amsmath,amssymb,amsthm,mathtools}
\usepackage[margin=2.5cm]{geometry}
\usepackage{hyperref}
\usepackage{booktabs}
\usepackage{enumitem}
\usepackage{array}

%% ---------- theorem environments ----------
\newtheorem{theorem}{Theorem}[section]
\newtheorem{proposition}[theorem]{Proposition}
\newtheorem{lemma}[theorem]{Lemma}
\newtheorem{corollary}[theorem]{Corollary}
\theoremstyle{definition}
\newtheorem{definition}[theorem]{Definition}
\newtheorem{example}[theorem]{Example}
\newtheorem{conjecture}[theorem]{Conjecture}
\newtheorem{observation}[theorem]{Observation}
\theoremstyle{remark}
\newtheorem{remark}[theorem]{Remark}

%% ---------- macros ----------
\DeclareMathOperator{\val}{v}
\newcommand{\N}{\mathbb{N}}
\newcommand{\Z}{\mathbb{Z}}
\newcommand{\Q}{\mathbb{Q}}
\newcommand{\R}{\mathbb{R}}
\newcommand{\F}{\mathbb{F}}
\newcommand{\Nodd}{\N_{\mathrm{odd}}}
\newcommand{\norm}[1]{\left\|#1\right\|}

%% ---------- metadata ----------
\title{\Large The $65/41$ Barrier:\\
Continued-Fraction Convergents of $\log_2 3$\\
Bound Worst-Case Descent Ratios\\[4pt]
\large in the Compressed Collatz Map}
\author{Lightman Chang\\
Independent Researcher\\
\texttt{lightman.chang@gmail.com}}
\date{}

%% =================================================================
\begin{document}
\maketitle

%% =================================================================
%% ABSTRACT
%% =================================================================
\begin{abstract}
For the compressed Collatz map $S(n) = (3n+1)/2^{\val_2(3n+1)}$ on the
positive odd integers, an orbit segment of length $T$ satisfies
$S^T(n) < n$ for all sufficiently large $n$ if and only if
$B_T = \sum_{i=1}^{T} \val_2(3 S^{i-1}(n) + 1)$ exceeds $T \log_2 3$.
The ratio $B_T / T$ is therefore a natural \emph{descent ratio} that
quantifies how close an orbit segment comes to balancing the
multiplicative growth by $3^T$ with the contraction by $2^{B_T}$.
We study the worst (smallest) value of this ratio across orbit segments
that descend.  Our main contribution (Theorem~\ref{thm:cf})
is a continued-fraction theorem: among all positive integers $T$ in the
range $1 \leq T \leq 93$, the quantity
$B(T) := \min\{B \in \Z : B/T > \log_2 3\}$ attains its minimum
$B(T)/T$ exactly at $T \in \{41, 82\}$ with value $65/41 = 1.58536\ldots$,
the sixth principal convergent of $\log_2 3$.  We complement this with a
computational study (exhaustive odd $n < 2^{21}$ and a structured sample
for $n < 2^{44}$) showing that the worst-case observed first-descent
ratio $B/T$ takes values in $\{65/41,\, 149/94,\, 233/147\}$~--- the
first three ``upper'' continued-fraction objects of $\log_2 3$~--- and
is at least $233/147$ in every bit-bin we tested.  As a rigorous
corollary (Corollary~\ref{cor:cycle}) we use the irrationality of
$\log_2 3$ to derive the cycle-element upper bound
$m < 2^B \cdot 3^P / \bigl(2(2^B - 3^P)\bigr)$ for any nontrivial cycle
of $S$.  All experimental code is publicly available.
\end{abstract}

\noindent\textbf{Keywords:}~Collatz map, $3x{+}1$ problem,
continued fractions, $\log_2 3$, descent ratio, computational number
theory.

\smallskip
\noindent\textbf{MSC 2020:}~Primary 11B37; Secondary 11A07, 11A55,
11J70, 11Y55, 11Y65, 37P99.

%% =================================================================
\section{Introduction}\label{sec:intro}
%% =================================================================

The Collatz map $T \colon \N \to \N$, defined by $T(n) = n/2$ for $n$
even and $T(n) = 3n + 1$ for $n$ odd, has resisted resolution of its
namesake conjecture (that every positive integer eventually reaches
$1$ under iteration of $T$) for nearly a century; see the surveys of
Lagarias~\cite{Lagarias1985,Lagarias2010} and the strongest known
density result of Tao~\cite{Tao2022}.  Throughout this paper we work
with the \emph{compressed} (or \emph{accelerated}) Collatz map
$S \colon \Nodd \to \Nodd$ defined by
\begin{equation}\label{eq:S-def}
  S(n) \;=\; \frac{3n + 1}{2^{\val_2(3n + 1)}},
\end{equation}
where $\val_2(m)$ is the $2$-adic valuation of $m$.  Iterating $S$
starting from an odd integer $n$ produces a sequence
$\val_1, \val_2, \val_3, \ldots$ of positive integer valuations, with
$\val_i = \val_2(3 S^{i-1}(n) + 1)$, and the iterate after $T$ steps is
\begin{equation}\label{eq:ST}
  S^T(n) \;=\; \frac{3^T n + C_T}{2^{B_T}},
  \qquad
  B_T \;=\; \sum_{i=1}^{T} \val_i,
\end{equation}
with a positive integer correction term $C_T$ depending on the
sequence $(\val_i)$ but not on $n$.

\paragraph{Descent ratios.}
For fixed $T$ and a fixed valuation sequence $(\val_1, \ldots, \val_T)$
realized by some orbit, the iterate $S^T(n) < n$ holds for all
$n > C_T / (2^{B_T} - 3^T)$ provided $2^{B_T} > 3^T$, that is,
provided
\begin{equation}\label{eq:descent}
  \frac{B_T}{T} \;>\; \log_2 3.
\end{equation}
The ratio $B_T / T$ is therefore the natural \emph{descent ratio} of
an orbit segment of length $T$: it is precisely the average $2$-adic
valuation along the segment, and the orbit segment can decrease only
when this average exceeds $\log_2 3 \approx 1.584962$.  The descent
condition~\eqref{eq:descent} is equivalent to the standard fact that
the empirical mean of $\val_2$ along an orbit must exceed $\log_2 3$
for the orbit to contract; see~\cite[\S 2]{Lagarias1985}.

\paragraph{The barrier.}
Since $\log_2 3$ is irrational (an elementary fact: $2^B = 3^P$ would
violate unique factorization in $\Z$), the descent ratio $B_T / T$
can never equal $\log_2 3$ for any positive integers $B_T, T$.  The minimum positive integer $B$ with
$B / T > \log_2 3$ is therefore
\[
  B(T) \;:=\; \lfloor T \log_2 3 \rfloor + 1.
\]
The ratio $B(T) / T$ is the smallest descent ratio achievable by an
orbit segment of length $T$ under the constraint
$B_T/T > \log_2 3$.  This deterministic minimum is governed by the
continued-fraction expansion of $\log_2 3$.  By the standard theory
\cite[Ch.~7]{HardyWright}, the principal convergents $h_k = p_k/q_k$
of $\log_2 3$ alternate above and below the irrational and are
``best one-sided rational approximations'' in the sense that any
fraction $p/q$ with $1 \leq q \leq q_{k+1} - 1$ and $p/q > \log_2 3$
satisfies $p/q \geq p_k/q_k$ when $h_k$ is an above-approximation.

\paragraph{The empirical observation.}
Computational experiments on compressed Collatz orbits (described in
detail in Section~\ref{sec:exp}) reveal the following striking
phenomenon: across all bit-shells $13 \leq k \leq 44$ that we tested
(with $\Omega_k = \Nodd \cap [2^{k-1},\, 2^k)$), the worst-case
first-descent ratio $B/T$ takes values only in the three-element set
$\{65/41,\, 149/94,\, 233/147\}$: it equals $65/41$ in most bit-shells,
with sporadic excursions to $149/94$ (at $k \in \{19, 21, 24, 27, 37,
42\}$ in our sample) and to $233/147$ (at $k = 44$).  All three
ratios are intimately tied to the continued-fraction expansion of
$\log_2 3$: $65/41$ is the sixth principal convergent, $149/94$ is the
next ``above'' semiconvergent, and $233/147$ is the next one after
that.

\paragraph{Our contribution.}
The main result of this paper is the continued-fraction barrier
theorem (Theorem~\ref{thm:cf}); the empirical and corollary
contributions are framed as supporting findings.

\begin{enumerate}[label=(\roman*),leftmargin=2em]
  \item \emph{Main result: Continued-fraction barrier theorem}
    (Theorem~\ref{thm:cf}, rigorous).
    For every positive integer $T$ with $1 \leq T \leq 93$ the
    quantity $B(T)/T$ satisfies $B(T)/T \geq 65/41$, with equality
    iff $T \in \{41, 82\}$.  The constant $65/41$ is the sixth
    principal convergent of $\log_2 3$, and the bound $T \leq 93$ is
    sharp: at $T = 94$ one has $B(94)/94 = 149/94 < 65/41$.
  \item \emph{Supporting empirical observation}
    (Conjecture~\ref{conj:65-41}, computational).
    Based on exhaustive enumeration for $n < 2^{21}$ and a structured
    sample for $n < 2^{44}$, we conjecture that for every odd $n$ of
    bit-length $k \geq 13$ the first-descent ratio satisfies
    $\rho_{\mathrm{des}}(n) \geq 233/147$, and that the worst-case
    bit-bin value lies in $\{65/41, 149/94, 233/147\}$.
  \item \emph{Cycle-element corollary}
    (Corollary~\ref{cor:cycle}, rigorous).
    For any nontrivial cycle of $S$ of period $P$ with cumulative
    valuation $B$, the smallest cycle element $m$ satisfies
    $m < 2^B \cdot 3^P / \bigl(2 (2^B - 3^P)\bigr)$, where
    $2^B - 3^P \geq 1$ by the irrationality of $\log_2 3$.  This is
    a recasting of a classical observation of Crandall (1978,
    \cite[Eq.~(8)]{Crandall1978}); we record it because the same
    arithmetic ($2^B \neq 3^P$) drives Theorem~\ref{thm:cf}.
\end{enumerate}

The structure of the paper places the rigorous theorems
(Theorem~\ref{thm:cf} in Section~\ref{sec:cf} and the cycle bound in
Section~\ref{sec:cycle}) on either side of the experimental
investigation (Sections~\ref{sec:exp} and~\ref{sec:conj}).  This
mirrors the actual epistemic state of the subject: the descent ratio
is forced from above by deterministic continued-fraction arithmetic,
the cycle existence problem is constrained from above by the
irrationality of $\log_2 3$, but the assertion that the empirical
worst case equals the
deterministic lower bound is at present an experimental observation
rather than a theorem.

\paragraph{What we do not claim.}
We make no claim of resolving the Collatz conjecture.  The empirical
match between worst-case observed descent ratios and the
continued-fraction lower bound is a precise statement about a
specific finite-range computation; extending it to all $n$ is a
genuine open problem, intimately tied to the equidistribution of
Collatz orbits modulo small powers of $2$.  We discuss the gap
carefully in Section~\ref{sec:disc}.

\paragraph{Relation to other papers in this series.}
The present paper is the third in a series of self-contained pieces
extracted from the same research notebook.  The first paper
\cite{ChangPaperA} constructs a three-state finite automaton for the
carry computation in $3n + 1$ and computes its spectral gap.  The
second paper \cite{ChangPaperB} proves an impossibility result for
finite-precision Lyapunov potentials.  The present paper isolates the
continued-fraction descent-ratio phenomenon and the associated
cycle-element bound; we use the closed-form shift identity from
\cite{ChangPaperA} as a black box in passing but the present paper
is logically self-contained.

\paragraph{Paper organization.}
Section~\ref{sec:prelim} fixes notation and recalls the basic
arithmetic of compressed Collatz iterates and the continued-fraction
expansion of $\log_2 3$.  Section~\ref{sec:cf} states and proves the
continued-fraction barrier theorem (Theorem~\ref{thm:cf}).
Section~\ref{sec:exp} describes the computational protocol and
reports the experimental results.  Section~\ref{sec:conj} formulates
the $65/41$ empirical conjecture (Conjecture~\ref{conj:65-41}) and
discusses its connection to the worst-case orbits we identified.
Section~\ref{sec:cycle} derives the cycle-element upper bound from
the irrationality of $\log_2 3$ as Corollary~\ref{cor:cycle}.
Section~\ref{sec:disc} discusses limitations, related work, and open
questions.

%% =================================================================
\section{Preliminaries}\label{sec:prelim}
%% =================================================================

Throughout, $\Nodd = \{1, 3, 5, \ldots\}$ denotes the positive odd
integers, and for $m \in \Z \setminus \{0\}$ we write $\val_2(m)$ for
the $2$-adic valuation of $m$.  We write
$\log_2 3 = 1.584962500721156\ldots$ for the base-$2$ logarithm of
$3$.

\subsection{Compressed Collatz iterates}

\begin{definition}[Compressed Collatz map]\label{def:S}
The \emph{compressed Collatz map} is the map
$S \colon \Nodd \to \Nodd$ defined by~\eqref{eq:S-def}.  For
$T \geq 1$ and $n \in \Nodd$ we write $S^T(n)$ for the $T$-th iterate
and define the \emph{valuation sequence} of $n$ at depth $T$ by
\[
  \val_i(n) \;=\; \val_2\bigl(3 S^{i-1}(n) + 1\bigr),
  \qquad 1 \leq i \leq T.
\]
The \emph{cumulative valuation} is $B_T(n) = \sum_{i=1}^{T} \val_i(n)$,
and the \emph{descent ratio} of the segment is $\rho_T(n) = B_T(n)/T$.
\end{definition}

\begin{lemma}[Closed form of compressed iterates]\label{lem:closed}
For every $n \in \Nodd$ and every $T \geq 1$,
\begin{equation}\label{eq:STexp}
  S^T(n) \;=\; \frac{3^T n + C_T(n)}{2^{B_T(n)}},
\end{equation}
where, with the convention $B_0(n) = 0$, the correction term
\[
  C_T(n) \;=\; \sum_{j=0}^{T-1} 3^{T-1-j}\, 2^{B_j(n)}
\]
is a positive integer depending only on the partial sums
$B_0(n), B_1(n), \ldots, B_{T-1}(n)$ (and not on $n$ apart from through
the $\val_i$).
\end{lemma}

\begin{proof}
Induction on $T$.  For $T = 1$ we have $S(n) = (3n + 1)/2^{\val_1}$
and $C_1 = 3^0 \cdot 2^{B_0} = 1$, agreeing with the definition.
Assume \eqref{eq:STexp} at depth $T$.  From
$S^{T+1}(n) = (3 S^T(n) + 1)/2^{\val_{T+1}}$ and the inductive
formula $S^T(n) = (3^T n + C_T)/2^{B_T}$ we obtain
\[
  S^{T+1}(n)
   = \frac{3^{T+1} n + 3 C_T + 2^{B_T}}{2^{B_T + \val_{T+1}}}
   = \frac{3^{T+1} n + C_{T+1}}{2^{B_{T+1}}},
\]
so $C_{T+1} = 3 C_T + 2^{B_T}$.  Substituting the inductive formula
and using $B_0 = 0$,
\[
  C_{T+1}
  = 3 \sum_{j=0}^{T-1} 3^{T-1-j} 2^{B_j} + 2^{B_T}
  = \sum_{j=0}^{T-1} 3^{T-j} 2^{B_j} + 2^{B_T}
  = \sum_{j=0}^{T} 3^{T-j} 2^{B_j}
  = \sum_{j=0}^{T} 3^{(T+1) - 1 - j} 2^{B_j},
\]
which is the claimed expression at depth $T+1$.
\end{proof}

\begin{corollary}[Descent criterion]\label{cor:descent}
For $n \in \Nodd$ and $T \geq 1$, the inequality $S^T(n) < n$ holds
if and only if $n (2^{B_T(n)} - 3^T) > C_T(n)$.  In particular, if
$\rho_T(n) > \log_2 3$, then $2^{B_T} > 3^T$ and the inequality
$S^T(n) < n$ holds for every
$n > C_T(n) / (2^{B_T(n)} - 3^T)$.
\end{corollary}

\begin{proof}
By Lemma~\ref{lem:closed},
$S^T(n) < n \iff (3^T n + C_T)/2^{B_T} < n
 \iff n (2^{B_T} - 3^T) > C_T$.  The condition
$\rho_T = B_T/T > \log_2 3$ is equivalent to
$2^{B_T} > 3^T$, hence $2^{B_T} - 3^T \geq 1$ since both are
positive integers.
\end{proof}

\subsection{Continued fractions of $\log_2 3$}

\begin{definition}[Convergents]\label{def:conv}
The simple continued-fraction expansion of $\log_2 3$ is
\[
  \log_2 3 \;=\; [a_0; a_1, a_2, \ldots]
  \;=\; [1;\, 1, 1, 2, 2, 3, 1, 5, 2, 23, \ldots],
\]
with partial quotients $a_0, a_1, a_2, \ldots$.  The principal
convergents $h_k = p_k / q_k$ are defined recursively by
\[
  p_{-1} = 1,\; p_{0} = a_0 = 1,\; p_k = a_k p_{k-1} + p_{k-2};
  \qquad
  q_{-1} = 0,\; q_0 = 1,\; q_k = a_k q_{k-1} + q_{k-2}.
\]
\end{definition}

We list the first few convergents in Table~\ref{tab:conv}.

\begin{table}[htbp]
\centering
\renewcommand{\arraystretch}{1.15}
\begin{tabular}{c c c r r}
\toprule
$k$ & $a_k$ & $h_k = p_k/q_k$ & decimal value & $h_k - \log_2 3$ \\
\midrule
$0$ & $1$ & $1/1$ & $1.000000$ & $-0.584962$ \\
$1$ & $1$ & $2/1$ & $2.000000$ & $+0.415037$ \\
$2$ & $1$ & $3/2$ & $1.500000$ & $-0.084962$ \\
$3$ & $2$ & $8/5$ & $1.600000$ & $+0.015037$ \\
$4$ & $2$ & $19/12$ & $1.583333$ & $-0.001629$ \\
$5$ & $3$ & $\mathbf{65/41}$ & $\mathbf{1.585366}$ &
  $\mathbf{+0.000403}$ \\
$6$ & $1$ & $84/53$ & $1.584906$ & $-0.000057$ \\
$7$ & $5$ & $485/306$ & $1.584967$ & $+0.0000048$ \\
$8$ & $2$ & $1054/665$ & $1.584962$ & $-0.000000095$ \\
\bottomrule
\end{tabular}
\caption{The first nine principal convergents of $\log_2 3$.  The
sixth convergent $h_5 = 65/41$ is the focus of the present paper:
it is the worst above-approximation of $\log_2 3$ in a sense made
precise by Theorem~\ref{thm:cf}.}
\label{tab:conv}
\end{table}

\begin{lemma}[Standard properties of convergents]\label{lem:cfprop}
With notation as above:
\begin{enumerate}[label=(\roman*),leftmargin=2em]
  \item $|p_k - q_k \log_2 3| < |p_{k-1} - q_{k-1} \log_2 3|$ for
    $k \geq 1$.
  \item $h_k$ alternates above and below $\log_2 3$ with
    $h_{2j+1} > \log_2 3 > h_{2j}$ for $j \geq 0$, and
    $h_5 = 65/41 > \log_2 3$.
  \item For any rational $p/q$ with $1 \leq q < q_{k+1}$ and
    $p/q > \log_2 3$, we have $p/q \geq h_k$ provided
    $h_k > \log_2 3$.
\end{enumerate}
\end{lemma}

\begin{proof}
These are standard consequences of the theory of continued fractions;
see, e.g., \cite[Theorems~163, 181, 182]{HardyWright}.
\end{proof}

\subsection{The descent denominator function}

We will repeatedly use the following deterministic minimum.

\begin{definition}[Descent denominator]\label{def:Bmin}
For each positive integer $T$, let
\[
  B(T) \;:=\; \lfloor T \log_2 3 \rfloor + 1,
\]
the smallest positive integer $B$ with $B / T > \log_2 3$.
\end{definition}

Since $\log_2 3$ is irrational, $T \log_2 3$ is never an integer, so
$B(T) = \lceil T \log_2 3 \rceil$ for every $T \geq 1$.

\begin{example}\label{ex:Bsmall}
We have $B(1) = 2$, $B(2) = 4$, $B(3) = 5$, $B(5) = 8$, $B(12) = 20$,
$B(17) = 27$, $B(29) = 46$, and $B(41) = 65$.  In particular
$B(41)/41 = 65/41 = 1.5853658\ldots$
\end{example}

%% =================================================================
\section{The Continued-Fraction Barrier Theorem}\label{sec:cf}
%% =================================================================

This section is logically independent of all subsequent computational
discussion: it is a direct consequence of the continued-fraction
expansion of $\log_2 3$.

\begin{theorem}[Continued-fraction barrier]\label{thm:cf}
For every integer $T$ with $1 \leq T \leq 93$,
\begin{equation}\label{eq:cf-bound}
  \frac{B(T)}{T} \;\geq\; \frac{65}{41} \;=\; 1.5853658536\ldots,
\end{equation}
with equality if and only if $T \in \{41, 82\}$.  Equivalently, the
sixth principal convergent $h_5 = 65/41$ of $\log_2 3$ realizes the
minimum of $B(T)/T$ over the range $1 \leq T \leq 93$.  The bound
$T \leq 93$ is sharp: at $T = 94$ one has
$B(94)/94 = 149/94 = 1.5851063\ldots < 65/41$.
\end{theorem}

The remainder of this section is the proof.  We argue in three steps:
first we identify $h_5$ as a best above-approximation in a precise
sense; second we use the standard alternation theorem to rule out
denominators below $41$; and third we directly verify the range
$42 \leq T \leq 93$.

\subsection{Best above-approximations of $\log_2 3$}

For an irrational $\alpha$ and a positive integer $T$, define
$\beta(T) := \lfloor T \alpha \rfloor + 1 - T \alpha = T (B(T)/T -
\alpha) > 0$, the ``positive defect'' of the smallest above-fraction
with denominator $T$.

\begin{lemma}[Defect minima at convergents]\label{lem:defect}
Let $\alpha = \log_2 3$.  For $T_0 \geq 1$, define
\[
  T^{*}(T_0) \;:=\; \operatorname*{arg\,min}_{1 \leq T \leq T_0} \frac{B(T)}{T}.
\]
Then $T^{*}(T_0)$ is the largest integer $T \leq T_0$ such that
$T \alpha - \lfloor T \alpha \rfloor$ is the maximum of
$\{j \alpha - \lfloor j \alpha \rfloor : 1 \leq j \leq T_0\}$, that
is, the integer $T$ for which the fractional part $\{T \alpha\}$ is
closest to $1$ from below.
\end{lemma}

\begin{proof}
We have $B(T)/T = (\lfloor T \alpha \rfloor + 1)/T = \alpha + (1 -
\{T \alpha\})/T$.  Since $\alpha$ is fixed, minimizing $B(T)/T$ over
$T \in [1, T_0]$ is the same as minimizing $(1 - \{T \alpha\})/T$.
For each fixed $T$, the numerator $1 - \{T \alpha\} > 0$, and we
seek the $T$ that makes this small (i.e.\ $\{T\alpha\}$ close to~$1$
from below) while $T$ is large.  The equivalent reformulation in the
statement is immediate.
\end{proof}

\begin{lemma}[Three-distance lemma applied to convergents]
\label{lem:bestabove}
The denominators $T \in \{1, 2, 3, \ldots\}$ at which $\{T \log_2 3\}$
attains a new maximum (in $[0, 1)$) are exactly the denominators of
those principal convergents $h_k = p_k/q_k$ of $\log_2 3$ that lie
\emph{above} $\log_2 3$, namely $h_1, h_3, h_5, h_7, \ldots$, together
with the corresponding ``upper semiconvergents''
\[
  \frac{p_{2j-1} + i \, p_{2j}}{q_{2j-1} + i \, q_{2j}},
  \quad 1 \leq i \leq a_{2j+1} - 1,
\]
that lie above $\log_2 3$ and lie strictly between $h_{2j-1}$ and
$h_{2j+1}$.  In particular, the first six denominators $T$ at which
the running maximum of $\{T \log_2 3\}$ strictly increases are
\[
  T \;\in\; \{1,\, 3,\, 5,\, 17,\, 29,\, 41\},
\]
with corresponding ratios $B(T)/T$ equal to
$2/1,\, 5/3,\, 8/5,\, 27/17,\, 46/29,\, 65/41$, respectively.
\end{lemma}

\begin{proof}
This is a textbook consequence of the three-distance theorem and the
theory of best one-sided rational approximations; see, e.g.,
\cite[\S 5--6]{Khinchin1964} or \cite[Ch.~6]{Bugeaud2004}.  The
denominators of best above-approximations of $\log_2 3$ are exactly
the upper semiconvergents (those rational expressions whose value
exceeds $\log_2 3$) along the sequence
$h_{1}, h_{3}, h_{5}, \ldots$  Direct computation from
Table~\ref{tab:conv} gives that the upper semiconvergents between
$h_3 = 8/5$ and $h_5 = 65/41$ are
\[
  \frac{8 + 1 \cdot 19}{5 + 1 \cdot 12} \;=\; \frac{27}{17},
  \qquad
  \frac{8 + 2 \cdot 19}{5 + 2 \cdot 12} \;=\; \frac{46}{29},
\]
since $a_5 = 3$ produces semiconvergents at $i = 1, 2$ before
reaching the next convergent at $i = 3$, which is
$h_5 = (8 + 3 \cdot 19)/(5 + 3 \cdot 12) = 65/41$.  Listing the
denominators $1, 3, 5, 17, 29, 41$ matches the claim.
\end{proof}

\begin{remark}\label{rem:T82}
The same fraction $65/41$ is also realized at $T = 82$ since
$B(82) = 130 = 2 \cdot 65$ and $130/82 = 65/41$.  This is consistent
with $\lfloor 82 \log_2 3 \rfloor + 1 = 129 + 1 = 130$.  Equality
$B(T)/T = 65/41$ in the range $1 \leq T \leq 93$ holds therefore at
$T = 41$ and $T = 82$, and at no other $T$ in this range.
\end{remark}

\subsection{Proof of Theorem~\ref{thm:cf}}

\begin{proof}[Proof of Theorem~\ref{thm:cf}]
We must show that $B(T)/T \geq 65/41$ for every $T \in [1, 93]$ and
that the bound is attained exactly at $T \in \{41, 82\}$.

\emph{Step 1: Reduction to record-breakers.}  By
Lemma~\ref{lem:bestabove}, the values of $T$ at which $B(T)/T$
strictly decreases (forming the running minimum) up to $T = 41$ are
exactly $T \in \{1, 3, 5, 17, 29, 41\}$.  At each subsequent integer
$T$ in $[1, 41]$ the value $B(T)/T$ is at least the previous record,
and the records take the values
$2/1, 5/3, 8/5, 27/17, 46/29, 65/41$, all of which are
$\geq 65/41$ by direct comparison
($65/41 = 1.585366$, $46/29 = 1.586207$, $27/17 = 1.588235$,
$8/5 = 1.6$, $5/3 = 1.667$, $2/1 = 2.0$).  The smallest among these
is $65/41$, attained at $T = 41$.

\emph{Step 2: $42 \leq T \leq 93$.}  We must show
$B(T)/T \geq 65/41$ for every such $T$, with equality only at $T = 82$.
By Lemma~\ref{lem:bestabove}, the running minimum of $B(T)/T$ is
attained, in the range $T \geq 1$, exactly at the denominators of
upper convergents/semiconvergents of $\log_2 3$.  The denominator
sequence at which this running minimum strictly decreases is
$1, 3, 5, 17, 29, 41, 94, \ldots$.  The next strict improvement
\emph{after} $T = 41$ is therefore at $T = 94$, with value $149/94$.
Concretely, $h_6 = 84/53 < \log_2 3$ is a below-approximation and
does not threaten the above-approximation $h_5$; the next
above-approximation is the upper semiconvergent
\[
  m_1 \;=\; \frac{p_5 + 1 \cdot p_6}{q_5 + 1 \cdot q_6}
       \;=\; \frac{65 + 84}{41 + 53}
       \;=\; \frac{149}{94}.
\]
Hence for every $T$ with $42 \leq T \leq 93$ and $T \neq 82$,
$B(T)/T > 65/41$ strictly, because by Lemma~\ref{lem:bestabove}
the next denominator at which the running minimum of $B(T)/T$
strictly improves below $65/41$ is $T = 94$, which lies outside
the range $T \leq 93$.  Equality $B(T)/T = 65/41$
in the range $42 \leq T \leq 93$ occurs \emph{only} at $T = 82$
(twice the denominator $41$, as in Remark~\ref{rem:T82}); this is
verified by direct calculation:
\[
  B(82) \;=\; \lfloor 82 \cdot \log_2 3 \rfloor + 1
   \;=\; \lfloor 129.9669\ldots \rfloor + 1
   \;=\; 130, \qquad 130/82 = 65/41.
\]

\emph{Step 3: Sharpness at $T = 94$.}  Direct computation gives
$94 \log_2 3 = 148.9866\ldots$, so $B(94) = 149$ and
$B(94)/94 = 149/94 = 1.585106\ldots < 65/41$.  Thus $T = 94$ is the
smallest positive integer $T$ for which $B(T)/T < 65/41$.

Combining Steps~1--3 yields the theorem.
\end{proof}

\begin{remark}[Why $h_5$ and not $h_3$ or $h_7$]
The reader may ask why the sixth convergent $h_5 = 65/41$, rather
than $h_3 = 8/5$ or $h_7 = 485/306$, is the empirically observed
value in Section~\ref{sec:exp}.  The answer is one of scale.  At
$T = 5$, the bound $B(T)/T \geq 8/5 = 1.6$ is comfortably above
$\log_2 3 = 1.585$; orbits of length $5$ have descent ratios
significantly above the worst-case asymptote.  By contrast, by the
time $T$ reaches $41$, the deterministic minimum $B(T)/T$ has
descended to $65/41 = 1.5854$, only $0.04\%$ above $\log_2 3$;
this is essentially the asymptotic value, and orbits of length
$\geq 41$ rarely achieve a substantially better ratio (the next
improvement only kicks in at $T = 94$ with $149/94$, a relative
improvement of just $0.016\%$).  Section~\ref{sec:exp} shows that
typical worst-case Collatz orbits in the bit-length range we
sampled hit lengths $T$ that are most often around $41$, $82$, or
small multiples thereof, locking in the $65/41$ value.
\end{remark}

\begin{remark}[Length of validity]
Theorem~\ref{thm:cf} holds for $T \leq 93$.  Beyond this range the
deterministic minimum drops further: from $T = 94$ to $T = 146$ the
bound is $149/94$, then $233/147$, then $317/200$, and so on, each
slightly closer to $\log_2 3$ from above.  The next principal
convergent above $\log_2 3$ is $h_7 = 485/306$, so this descent
proceeds in small steps.  The empirical study in
Section~\ref{sec:exp} works at scales (orbit lengths $T$ up to a few
hundred for $k \leq 44$) where these tighter bounds are theoretically
allowed but, as we shall see, only $149/94$ ever \emph{occurs} in
our data, never $233/147$ or anything beyond.
\end{remark}

%% =================================================================
\section{Computational Experiments}\label{sec:exp}
%% =================================================================

This section describes the computational protocol used to investigate
worst-case descent ratios of compressed Collatz orbits and reports
the results.  The code is publicly available at
\url{https://github.com/weiqi-kids/collatz-research-report}.

\subsection{Setup and definitions}

\begin{definition}[Length of first descent]\label{def:Tdes}
For an odd integer $n \geq 3$, the \emph{length of first descent}
$T_{\mathrm{des}}(n)$ is the smallest $T \geq 1$ such that
$S^T(n) < n$, when such a $T$ exists.  We set
$T_{\mathrm{des}}(n) = \infty$ otherwise.
\end{definition}

By Lemma~\ref{lem:closed}, for every $n \geq 3$ that eventually
descends below itself (which, by exhaustive verification of the
Collatz conjecture for $n \leq 2^{68}$~\cite{BarinaCheckup}, includes
every $n$ in our search range), the length $T_{\mathrm{des}}(n)$ is
finite.  We then define
\[
  B_{\mathrm{des}}(n) \;:=\; B_{T_{\mathrm{des}}(n)}(n),
  \qquad
  \rho_{\mathrm{des}}(n) \;:=\;
    B_{\mathrm{des}}(n) / T_{\mathrm{des}}(n).
\]

\begin{definition}[Worst-case descent ratio in a bit-bin]
\label{def:worstk}
For $k \geq 2$, let
$\Omega_k = \Nodd \cap [2^{k-1},\, 2^k)$ denote the odd integers
of \emph{exactly} $k$ bits (i.e.\ the bit-shell of length $k$); for
$k = 2$ we set $\Omega_2 = \{3\}$.  The \emph{worst-case descent
ratio} in bit-bin $k$ is
\[
  \rho^{*}(k) \;:=\;
  \min\bigl\{\rho_{\mathrm{des}}(n) : n \in \Omega_k\bigr\}.
\]
By Theorem~\ref{thm:cf} (and the trivial extension to $T \geq 94$),
$\rho^{*}(k) > \log_2 3$ holds for every finite $k$.  We use
bit-shells (rather than cumulative ranges) so that distinct bit-bins
sample disjoint sets of integers; in particular, $\rho^{*}(k)$ is not
constrained to be monotone in $k$.
\end{definition}

\subsection{Algorithm}

We use the simple procedure of Table~\ref{tab:alg-descent} to compute
$T_{\mathrm{des}}(n)$, $B_{\mathrm{des}}(n)$, and the corresponding
descent ratio.

\begin{table}[htbp]
\centering
\begin{tabular}{|p{0.92\linewidth}|}
\hline
\textbf{Algorithm 1.} \texttt{first\_descent\_ratio}$(n)$.\\
\hline
\textbf{Input:} odd integer $n \geq 3$.\\
\textbf{Output:} pair $(T_{\mathrm{des}}, B_{\mathrm{des}})$ such that
$S^{T_{\mathrm{des}}}(n) < n$ and $T_{\mathrm{des}}$ is minimal with
this property.\\[0.4em]
\quad 1. $m \leftarrow n$;\, $T \leftarrow 0$;\, $B \leftarrow 0$.\\
\quad 2. \textbf{while} $m \geq n$ \textbf{do}\\
\quad\quad 2a. $u \leftarrow 3m + 1$;\, $v \leftarrow \val_2(u)$.\\
\quad\quad 2b. $m \leftarrow u / 2^v$;\, $T \leftarrow T + 1$;\,
$B \leftarrow B + v$.\\
\quad 3. \textbf{return} $(T, B)$.\\
\hline
\end{tabular}
\caption{Algorithm \texttt{first\_descent\_ratio}, formatted as a
table for typographic convenience; we will refer to it as
``Algorithm~1''.}
\label{tab:alg-descent}
\end{table}

The algorithm terminates for every $n$ in our search range, since
the compressed Collatz iterate has been verified down to $1$ for all
$n \leq 2^{68}$~\cite{BarinaCheckup}, and a $T$-step descent below
$n$ occurs in a finite prefix of the orbit.

\subsection{Experimental design}

We carry out three increasingly large experiments.

\paragraph{Experiment E1 (exhaustive enumeration).}
For each $k \in \{5, 6, \ldots, 21\}$, we compute
$T_{\mathrm{des}}(n)$, $B_{\mathrm{des}}(n)$ and
$\rho_{\mathrm{des}}(n)$ for every odd $n \in \Omega_k$, and record
$\rho^{*}(k)$ as well as the integer pair $(T, B)$ that realizes it.
The total work is on the order of $2^{20} \approx 1.0 \times 10^6$
descent computations, each running in microseconds; total runtime is
under a minute on a single core.

\paragraph{Experiment E2 (Mersenne and trailing-$1$ orbits).}
For each $k \in \{5, 10, 15, 20, 25, 30, 35, 40, 44\}$, we examine
the Mersenne orbit starting at $n = 2^k - 1$ and the orbits starting
at $n = 2^k m' - 1$ for small odd $m'$.  These are theoretically
known to have long initial bad-step segments (by the closed-form
shift identity $S^{t-1}(2^t m' - 1) = 2 \cdot 3^{t-1} m' - 1$ from
\cite{ChangPaperA}, which we use as a black box) and are therefore
candidate sources of worst-case descent ratios.

\paragraph{Experiment E3 (structured sampling for $22 \leq k \leq 44$).}
For each $k$ in this range, we compute $\rho_{\mathrm{des}}(n)$ for a
mixed sample drawn from four sources, totaling between $10^4$ and
$10^5$ samples per $k$:
\begin{enumerate}[label=(\alph*),leftmargin=2em]
  \item the Mersenne number $2^k - 1$ and the integers
    $2^k m' - 1$ for odd $m' \in \{1, 3, 5, \ldots, 999\}$
    (\emph{trailing-$1$ patterns}, contributing the long bad-step
    segments);
  \item the integer that achieved $\rho^{*}(k - 1)$ in the previous
    bit-bin together with its small odd multiples
    (\emph{previous-round expansion});
  \item $10^4$ pseudo-random odd integers drawn uniformly from
    $\Omega_k \setminus \Omega_{k-1}$ (\emph{random});
  \item all integers congruent to $-1$ modulo $2^{k - 6}$ within
    $\Omega_k$ (\emph{near-Mersenne mosaic}).
\end{enumerate}
This sampling strategy is designed to capture orbits that empirically
produce small descent ratios, complementing the exhaustive enumeration
at lower $k$.  We do not claim that the sample exhausts all worst-case
candidates; this is precisely the empirical character of the
observation discussed in Section~\ref{sec:conj}.

\paragraph{Environment.}
The experiments were run on a single workstation (Intel x86-64,
\texttt{Python 3.11}, native big-integer arithmetic).  Total wall-clock
time is dominated by E3 and amounts to under one CPU-hour.

\subsection{Results}

\paragraph{Experiment E1.}
Table~\ref{tab:E1} reports $\rho^{*}(k)$ for $k = 5, \ldots, 18$ in
fully-enumerated form.  The realizing $(T, B)$ pair is given as a
fraction $B/T$ in lowest terms, alongside the matching
continued-fraction object.

\begin{table}[htbp]
\centering
\renewcommand{\arraystretch}{1.15}
\begin{tabular}{r r r l l l}
\toprule
$k$ & $\#\Omega_k$ & worst-case $n$ & $(T,B)$ & $B/T$ (lowest terms)
  & continued-fraction interpretation \\
\midrule
$5$  & $8$        & $27$       & $(37,59)$  & $59/37$  & non-canonical small-$T$ \\
$6$  & $16$       & $63$       & $(34,54)$  & $27/17$  & upper semiconvergent \\
$7$  & $32$       & $71$       & $(32,51)$  & $51/32$  & non-canonical small-$T$ \\
$8$  & $64$       & $175$      & $(5,8)$    & $8/5$    & $h_3$ \\
$9$  & $128$      & $295$      & $(5,8)$    & $8/5$    & $h_3$ \\
$10$ & $256$      & $795$      & $(17,27)$  & $27/17$  & upper semiconvergent \\
$11$ & $512$      & $1051$     & $(15,24)$  & $8/5$    & $h_3$ \\
$12$ & $1024$     & $3431$     & $(29,46)$  & $46/29$  & upper semiconvergent \\
$13$ & $2048$     & $4591$     & $(41,65)$  & $65/41$  & $h_5$ \\
$14$ & $4096$     & $8255$     & $(17,27)$  & $27/17$  & upper semiconvergent (anomaly, see text) \\
$15$ & $8192$     & $26623$    & $(41,65)$  & $65/41$  & $h_5$ \\
$16$ & $16384$    & $35295$    & $(41,65)$  & $65/41$  & $h_5$ \\
$17$ & $32768$    & $71451$    & $(41,65)$  & $65/41$  & $h_5$ \\
$18$ & $65536$    & $154079$   & $(41,65)$  & $65/41$  & $h_5$ \\
$19$ & $131072$   & $432923$   & $(94,149)$ & $149/94$ & next semiconvergent \\
$20$ & $262144$   & $525159$   & $(41,65)$  & $65/41$  & $h_5$ \\
$21$ & $524288$   & $1130495$  & $(94,149)$ & $149/94$ & next semiconvergent \\
\bottomrule
\end{tabular}
\caption{Experiment E1: exhaustive worst-case first-descent ratios
$\rho^{*}(k) = \min_{n \in \Omega_k} \rho_{\mathrm{des}}(n)$ over the
bit-shell $\Omega_k = \Nodd \cap [2^{k-1},\, 2^k)$ for $5 \leq k \leq
21$.  ``$h_3 = 8/5$'' and ``$h_5 = 65/41$'' denote principal
convergents of $\log_2 3$; ``upper semiconvergent'' denotes a
best-above approximation between two principal convergents (see
Lemma~\ref{lem:bestabove}); ``non-canonical small-$T$'' denotes a
ratio realized in the small-$T$ regime that is not itself an
upper convergent or upper semiconvergent (e.g.\ $59/37 = 1.5946$,
$51/32 = 1.5938$).  Because the $\Omega_k$ are disjoint shells,
$\rho^{*}(k)$ is not constrained to be monotone in $k$: at $k = 14$,
for example, no orbit attains the length-$41$ first-descent value
$65/41$, so the worst case rolls back to $27/17$.  The exhaustive
enumeration was carried out for all $5 \leq k \leq 21$ via the
script \texttt{compute\_descent.py} (see Reproducibility,
Section~\ref{sec:disc}).}
\label{tab:E1}
\end{table}

\paragraph{Experiment E2.}
For Mersenne starting points $n = 2^k - 1$ we observed (with the
notation of Definition~\ref{def:S}):

\begin{table}[htbp]
\centering
\begin{tabular}{r r r r l r}
\toprule
$k$ & $n = 2^k - 1$ & $T_{\mathrm{des}}$ & $B_{\mathrm{des}}$
  & $B/T$ (reduced) & decimal \\
\midrule
$5$  & $31$            & $35$  & $56$  & $8/5$    & $1.6000$ \\
$10$ & $1023$          & $11$  & $20$  & $20/11$  & $1.8182$ \\
$15$ & $32767$         & $32$  & $52$  & $13/8$   & $1.6250$ \\
$20$ & $1048575$       & $34$  & $54$  & $27/17$  & $1.5882$ \\
$25$ & $33554431$      & $92$  & $150$ & $75/46$  & $1.6304$ \\
$30$ & $1073741823$    & $76$  & $122$ & $61/38$  & $1.6053$ \\
$41$ & $2199023255551$ & $75$  & $120$ & $8/5$    & $1.6000$ \\
$44$ & $17592186044415$& $128$ & $204$ & $51/32$  & $1.5938$ \\
\bottomrule
\end{tabular}
\caption{Experiment E2: first-descent statistics for the Mersenne
starting points $n = 2^k - 1$.  Note that Mersenne starts do not, in
general, realize the worst-case bit-bin ratio: their first-descent
ratios are typically substantially above the bin minimum (compare
with Table~\ref{tab:E1}, e.g.\ at $k = 20$ the bin worst case is
$65/41 = 1.5854$ at $n = 525159$, whereas the Mersenne starts at
$27/17 = 1.5882$).  The reason is that the Mersenne orbit
$2^k - 1 \to 2 \cdot 3^{k-1} - 1 \to \cdots$ proceeds through a
prescribed shift identity that produces $\val_2 = 1$ for many
consecutive steps, but does not in itself force the optimal
$\rho_{\mathrm{des}}$ at any specific $T$.}
\label{tab:E2}
\end{table}

\begin{remark}\label{rem:mersenne-not-saturating}
The Mersenne orbit at $k = 41$ has $T_{\mathrm{des}} = 75$, $B = 120$
and ratio $B/T = 8/5$, not $65/41$.  The integer that realizes
$\rho^{*}(20) = 65/41$ in Table~\ref{tab:E1} (at the smallest
bit-length where the value $65/41$ first appears together with a
sample size large enough for the optimum to occur) is
$n = 525159$ at $(T, B) = (41, 65)$.  Hence the empirical worst case
is realized by integers other than the Mersenne starting points; the
Mersenne family is theoretically interesting because of the
shift-identity structure of \cite{ChangPaperA} but is not the
extremal family for the descent ratio statistic.
\end{remark}

\paragraph{Experiment E3.}
For $22 \leq k \leq 44$ we report the worst-case observed descent
ratio $\rho^{*}_{\mathrm{obs}}(k)$ in Table~\ref{tab:E3}.

\begin{table}[htbp]
\centering
\begin{tabular}{r l l}
\toprule
bit-bin $k$ & $\rho^{*}_{\mathrm{obs}}(k)$ & continued-fraction interpretation \\
\midrule
$22$ & $65/41$ & $h_5$ \\
$23$ & $65/41$ & $h_5$ \\
$24$ & $149/94$ & next semiconvergent \\
$25$ & $65/41$ & $h_5$ \\
$26$ & $65/41$ & $h_5$ \\
$27$ & $149/94$ & next semiconvergent \\
$28$ to $36$ & $65/41$ in every $k$ & $h_5$ \\
$37$ & $149/94$ & next semiconvergent \\
$38$ to $41$ & $65/41$ in every $k$ & $h_5$ \\
$42$ & $149/94$ & next semiconvergent \\
$43$ & $65/41$ & $h_5$ \\
$44$ & $233/147$ & next semiconvergent \\
\bottomrule
\end{tabular}
\caption{Experiment E3: worst-case observed first-descent ratios in
bit-bins $22 \leq k \leq 44$ over the structured sample of
$\sim 4.5 \times 10^4$ orbits per bin (combining trailing-$1$
patterns, near-Mersenne mosaics, prev-round expansions, and three
random seeds of $1.5 \times 10^4$ uniform odd integers each).  All
observed values lie in $\{65/41,\, 149/94,\, 233/147\}$, the first
three ``upper'' rational approximations of $\log_2 3$ from above
(see Lemma~\ref{lem:bestabove}).  Excursions to $149/94$ at
$k \in \{24, 27, 37, 42\}$ and to $233/147$ at $k = 44$ are sporadic
and depend on the sample; an exhaustive enumeration at these $k$
values is currently infeasible at single-workstation scale.}
\label{tab:E3}
\end{table}

\paragraph{Distributional remarks.}
For the worst-case orbit at $k = 18$ ($n = 154{,}079$, with
$T_{\mathrm{des}} = 41$ and $B_{\mathrm{des}} = 65$), the
distribution of valuations along the descent is
\begin{equation}\label{eq:dist}
  \#\{i : \val_i = 1\} = 28,
  \quad
  \#\{i : \val_i = 2\} = 5,
  \quad
  \#\{i : \val_i = 3\} = 7,
  \quad
  \#\{i : \val_i = 6\} = 1.
\end{equation}
Hence $B = 28 + 10 + 21 + 6 = 65$ and $T = 41$, so
$\rho = 65/41$ exactly.  We emphasize that this distribution is
heavily skewed toward $\val = 1$ (the minimum possible value), and a
single occurrence of $\val = 6$ in the segment is what allows the
ratio to remain at the boundary value $65/41$.

%% =================================================================
\section{The $65/41$ Empirical Conjecture}\label{sec:conj}
%% =================================================================

We now formalize the empirical observation of Section~\ref{sec:exp}
into an explicit conjecture.  We emphasize that this is a
\emph{conjecture}, not a theorem: the supporting evidence is
computational and confined to bit-lengths $k \leq 44$.

\begin{conjecture}[The descent-ratio barrier conjecture]\label{conj:65-41}
There exists a threshold $k_0$ (with $k_0 = 13$ consistent with the
data of Table~\ref{tab:E1}) such that for every odd integer $n$ of
bit-length at least $k_0$, the first-descent ratio satisfies
\begin{equation}\label{eq:conj-65-41}
  \rho_{\mathrm{des}}(n) \;=\;
  \frac{B_{\mathrm{des}}(n)}{T_{\mathrm{des}}(n)}
  \;\geq\; \frac{233}{147}.
\end{equation}
Moreover, the worst-case (smallest) value of $\rho_{\mathrm{des}}(n)$
over $n$ in each bit-bin $\Omega_k$ with $k \geq k_0$ lies in the
finite set $\{65/41,\, 149/94,\, 233/147,\, \ldots\}$ of values
$B(T)/T$ at successive upper continued-fraction objects of $\log_2 3$
(principal convergents and upper semiconvergents), and at the bit-bin
scales accessible to our experiment (up to $k = 44$) only the first
three of these values, $\{65/41,\, 149/94,\, 233/147\}$, are
realized.
\end{conjecture}

\begin{remark}\label{rem:asymptotic}
Conjecture~\ref{conj:65-41} is the natural extrapolation of
Theorem~\ref{thm:cf} to first descents at orbit lengths beyond
$T = 93$.  The deterministic minimum of $B(T)/T$ continues to
decrease, in well-defined steps, at the denominators of further
upper convergents and upper semiconvergents
(Lemma~\ref{lem:bestabove}): the values $65/41 < 149/94 < 233/147 <
\cdots$ are the values of $B(T)/T$ at $T \in \{41, 94, 147, \ldots\}$
respectively, each strictly smaller than the previous and each
strictly larger than $\log_2 3$.  Conjecture~\ref{conj:65-41} asserts
that, in our sampled bit-length range, only the first \emph{three}
such values appear as bit-bin worst cases; deeper search at larger
$k$ may reveal further values such as $317/200$ and $485/306 = h_7$,
but at the same time the conjecture posits that the floor never
slides arbitrarily close to $\log_2 3$ at any \emph{single}
$n$.  The small-$T$ (i.e.\ $k \leq 12$) regime in
Table~\ref{tab:E1} realizes other ratios such as $59/37$, $51/32$,
$8/5$, $27/17$, $46/29$ that are not in this set; the threshold
$k \geq 13$ is the natural boundary of validity of the conjecture
as stated.
\end{remark}

\begin{remark}\label{rem:weaker}
A weaker form of Conjecture~\ref{conj:65-41} that may be more
tractable is:
\[
  \liminf_{n \to \infty} \rho_{\mathrm{des}}(n) \;\geq\; \log_2 3.
\]
This weaker form is implied by the Collatz conjecture itself
(every orbit eventually reaches~$1$, hence the average valuation
exceeds $\log_2 3$ over sufficiently long segments) but is not known
unconditionally; Tao's almost-everywhere bound~\cite{Tao2022} gives
the inequality almost surely with respect to logarithmic density.
The strong form $\geq 149/94$ in Conjecture~\ref{conj:65-41} is what
the data of Section~\ref{sec:exp} actually support.
\end{remark}

\begin{observation}[Worst-case orbit shape, $k \geq 13$]\label{obs:shape}
For every bit-bin $k$ with $13 \leq k \leq 44$ in
Section~\ref{sec:exp}, the length of first descent $T_{\mathrm{des}}$
realizing $\rho^{*}_{\mathrm{obs}}(k)$ takes one of the values
$T_{\mathrm{des}} \in \{17, 41, 94, 147\}$: $T = 41$ realizes the
ratio $65/41$, $T = 17$ realizes the same numerical value $27/17$
(at $k = 14$ only, where no length-$41$ orbit attains the convergent
value within the shell), $T = 94$ realizes $149/94$, and $T = 147$
realizes $233/147$ (at $k = 44$ only).  No worst-case orbit
identified in this range realized $T_{\mathrm{des}}$ strictly
greater than $147$.  The small-$T$ regime $k \leq 12$ (see
Table~\ref{tab:E1}) realizes worst-case denominators
$T_{\mathrm{des}} \in \{5, 15, 17, 29, 32, 34, 37\}$ at upper
semiconvergents below $h_5$ or in the non-canonical small-$T$ range.
\end{observation}

\begin{remark}[Connection to length statistics]
Observation~\ref{obs:shape} relates Conjecture~\ref{conj:65-41} to
the statistics of $T_{\mathrm{des}}$.  If a future result established
that for sufficiently large $n$, the length $T_{\mathrm{des}}(n)$
either lies in $\{1, 2, \ldots, 93\}$ or grows large enough that
the empirical valuation mean concentrates above $\log_2 3$ (in a
quantitative sense), then Conjecture~\ref{conj:65-41} would follow.
The first regime is the regime governed by Theorem~\ref{thm:cf};
the second regime is governed by the equidistribution of orbits
modulo small powers of $2$, which is the content of the questions
left open by~\cite{Tao2022}.
\end{remark}

%% =================================================================
\section{Cycle-Element Upper Bound}\label{sec:cycle}
%% =================================================================

We close with a rigorous corollary of the basic descent equation
combined with the irrationality of $\log_2 3$.  This section is
logically independent of the experimental Sections~\ref{sec:exp}
and~\ref{sec:conj}.

\subsection{The cycle equation}

\begin{definition}[Cycle of $S$]\label{def:cycle}
A \emph{cycle} of the compressed Collatz map of period $P \geq 1$ is
a finite sequence $(m_0, m_1, \ldots, m_{P-1})$ of odd integers with
$m_{i+1} = S(m_i)$ for $0 \leq i \leq P - 2$ and $S(m_{P-1}) = m_0$.
The cycle is \emph{nontrivial} if $\{m_0, \ldots, m_{P-1}\} \neq
\{1\}$ (the trivial cycle is the fixed point $S(1) = 1$ of period
$P = 1$).
\end{definition}

\begin{lemma}[Cycle equation]\label{lem:cycleeq}
Let $(m_0, m_1, \ldots, m_{P-1})$ be a cycle of $S$ of period $P$,
let $\val_i = \val_2(3 m_{i-1} + 1)$ for $1 \leq i \leq P$, and let
$B = \sum_{i=1}^{P} \val_i$.  Then
\begin{equation}\label{eq:cycle}
  m_0 \bigl(2^B - 3^P\bigr) \;=\; \sum_{j=0}^{P-1}
    3^{P - 1 - j} \, 2^{B_j},
\end{equation}
where $B_0 = 0$ and $B_j = \sum_{i=1}^{j} \val_i$ for $j \geq 1$.
In particular, the right-hand side is a positive integer.
\end{lemma}

\begin{proof}
Apply Lemma~\ref{lem:closed} with $T = P$ to the orbit starting at
$m_0$.  The condition $S^P(m_0) = m_0$ becomes
$(3^P m_0 + C_P)/2^B = m_0$, i.e.\ $m_0(2^B - 3^P) = C_P$, with
$C_P = \sum_{j=0}^{P-1} 3^{P-1-j} 2^{B_j}$ as in
Lemma~\ref{lem:closed} (with $B_0 = 0$).  Since each term
$3^{P-1-j} 2^{B_j}$ is a positive integer, the sum is a positive
integer.
\end{proof}

\subsection{Transcendence input}

\begin{lemma}[Irrationality of $\log_2 3$]\label{lem:gs}
For all positive integers $B$ and $P$, $2^B \neq 3^P$.
\end{lemma}

\begin{proof}
By unique prime factorization in $\Z$, $2^B$ has only the prime $2$
in its factorization while $3^P$ has only the prime $3$.  Since these
factorizations differ for any positive $B, P$, we have $2^B \neq 3^P$.
\end{proof}

\begin{remark}\label{rem:gs-elementary}
The conclusion $2^B \neq 3^P$ is equivalent to the irrationality of
$\log_2 3$ (any equality $2^B = 3^P$ would yield $\log_2 3 = B/P$).
The full Gelfond--Schneider theorem~\cite{Gelfond1934,Schneider1934}
yields the stronger fact that $\log_2 3$ is transcendental and, more
quantitatively, allows one to give effective lower bounds for
$|2^B - 3^P|$ via Baker's theorem~\cite{Baker1990}.  We do not pursue
those quantitative bounds here; only the elementary inequality
$2^B \neq 3^P$ is needed for Corollary~\ref{cor:cycle}.
\end{remark}

\subsection{The cycle-element bound}

\begin{corollary}[Cycle-element upper bound]\label{cor:cycle}
Let $(m_0, m_1, \ldots, m_{P-1})$ be a nontrivial cycle of the
compressed Collatz map of period $P$ with cumulative valuation
$B = \sum_{i=1}^{P} \val_i$.  Then $2^B - 3^P \geq 1$, and the
smallest element $m_{\min} = \min_i m_i$ of the cycle satisfies
\begin{equation}\label{eq:cyclebound}
  m_{\min} \;<\; \frac{2^B \cdot 3^P}{2 \, (2^B - 3^P)}.
\end{equation}
In particular, $B > P \log_2 3$, so $B/P > \log_2 3$ for every
nontrivial cycle.
\end{corollary}

\begin{proof}
Existence of a cycle implies existence of $m_0 > 0$ satisfying
\eqref{eq:cycle}.  The right-hand side of \eqref{eq:cycle} is positive
and $m_0 > 0$, so $2^B - 3^P > 0$.  By Lemma~\ref{lem:gs} the
quantities $2^B$ and $3^P$ are distinct positive integers, so
$2^B - 3^P \geq 1$.

To bound $m_0$ from above, observe from \eqref{eq:cycle} that
each $B_j \leq B - 1$ for $0 \leq j \leq P - 1$ (since
$B_j + \sum_{i=j+1}^{P} \val_i = B$ and each $\val_i \geq 1$), hence
$2^{B_j} \leq 2^{B-1}$.  Therefore
\begin{align*}
  m_0 (2^B - 3^P)
  &= \sum_{j=0}^{P-1} 3^{P-1-j} \, 2^{B_j}
  \;\leq\; 2^{B-1} \sum_{j=0}^{P-1} 3^{P-1-j}
  \;=\; 2^{B-1} \cdot \frac{3^P - 1}{3 - 1} \\
  &< \frac{2^{B-1} \cdot 3^P}{2}
   \;=\; \frac{2^B \cdot 3^P}{4}.
\end{align*}
Dividing by $2^B - 3^P \geq 1$ yields the stronger inequality
$m_0 < 2^B \cdot 3^P / (4(2^B - 3^P))$, which in particular implies
the bound stated in~\eqref{eq:cyclebound}.

The same argument applies starting from any cyclic rotation of the
sequence (changing the starting point $m_0$ to $m_i$); choosing
$i$ such that $m_i = m_{\min}$ proves the bound for the smallest
element.

Finally, $2^B - 3^P \geq 1 > 0$ gives $2^B > 3^P$, so
$B > P \log_2 3$, i.e.\ $B/P > \log_2 3$.
\end{proof}

\begin{remark}[Numerical instance]\label{rem:cyclenum}
Setting $(B, P) = (65, 41)$ in Corollary~\ref{cor:cycle} gives
\[
  2^{65} - 3^{41}
  \;=\; 36{,}893{,}488{,}147{,}419{,}103{,}232
      - 36{,}472{,}996{,}377{,}170{,}786{,}403
  \;=\; 420{,}491{,}770{,}248{,}316{,}829,
\]
hence by the sharper $1/4$ form of Corollary~\ref{cor:cycle},
$m_{\min} < 2^{65} \cdot 3^{41} /
(4 \cdot 420{,}491{,}770{,}248{,}316{,}829) \approx
8.0 \times 10^{20}$.  Since the Collatz conjecture has been verified
for $n \leq 2^{68} \approx 2.95 \times 10^{20}$~\cite{BarinaCheckup},
the bound for $(B, P) = (65, 41)$ is not strong
enough to rule out a cycle of this $(B, P)$ shape unconditionally.
This is consistent with Eliahou's lower bound for nontrivial cycle
periods~\cite{Eliahou1993}, which forces $P \geq 17{,}087{,}915$, a
constraint orders of magnitude beyond the $(B, P) = (65, 41)$ regime
considered here.
\end{remark}

\begin{remark}[Why the bound is not enough]\label{rem:gap}
The bound in Corollary~\ref{cor:cycle} grows roughly as $3^P / P$
(for $B \approx P \log_2 3$, the denominator $2^B - 3^P$ can be
much smaller than $3^P$).  Combined with Eliahou's lower bound on
$P$, the resulting bound on $m_{\min}$ is enormous, far exceeding
the verified range~\cite{BarinaCheckup}.  Stronger bounds require a
\emph{linear forms in logarithms} approach (Baker--style theorems;
see~\cite{Baker1990}) on $|B \log 2 - P \log 3|$, which give
quantitative lower bounds for $|2^B - 3^P|$ (and hence upper bounds
on $m_{\min}$) that decay only polynomially in $P$, still leaving a
large gap.  Closing this gap is one of the standard open problems in
the cycle-elimination literature; see~\cite{Steiner1977} for a
classical attack restricted to ``$1$-cycles''.
\end{remark}

%% =================================================================
\section{Discussion}\label{sec:disc}
%% =================================================================

\paragraph{The role of $65/41$ in the descent ratio landscape.}
Theorem~\ref{thm:cf} pinpoints $65/41$ as the deterministic minimum
of $B(T)/T$ in the range $T \leq 93$, and Conjecture~\ref{conj:65-41}
asserts the (much stronger) empirical fact that no first-descent
orbit segment of any length attains a ratio strictly between
$\log_2 3$ and $65/41$.  The first statement is a fact about the
continued fractions of $\log_2 3$; the second statement involves
specific Collatz orbits and is, at the present level of theory,
inaccessible.

\paragraph{What the empirical match does and does not establish.}
The empirical match between worst-case observed descent ratios and
the deterministic continued-fraction minimum at $T \leq 93$ is
striking but constitutes \emph{evidence}, not proof.  Specifically,
the data of Section~\ref{sec:exp} establish that for every odd
$n < 2^{21}$ of bit-length $k \geq 13$, and for every $n$ in our
structured sample for $2^{21} \leq n < 2^{44}$, the inequality
$\rho_{\mathrm{des}}(n) \geq 233/147 = 1.585034\ldots$ holds; the
small-$T$ regime $k \leq 12$ realizes other ratios above this
threshold (Table~\ref{tab:E1}).  They do not establish that the
$\geq 233/147$ inequality persists for all $n$.  Two cautionary
observations are in order.

First, even granting Conjecture~\ref{conj:65-41}, this would not
directly imply the Collatz conjecture: a divergent orbit could in
principle exist while its first-descent segments still have
ratio $\geq 65/41$ (the orbit would simply contain only finitely
many descents below the starting value).  The conjecture rules out
descent ratios approaching $\log_2 3$ from above too quickly, but it
does not rule out orbits that fail to descend at all.  Equivalently,
Conjecture~\ref{conj:65-41} is a quantitative refinement of the
``Lemma D'' descent statement (every orbit has \emph{some} descent
window) but does not prove that statement; the latter remains
equivalent to the Collatz conjecture itself.

Second, the empirical sample for $22 \leq k \leq 44$ in Experiment
E3 is structured but not exhaustive.  The structured sample is
designed to capture orbits that should produce small descent ratios
(trailing-$1$ patterns, near-Mersenne, expansions of previously
worst-case orbits), and the consistency of the result over $23$
distinct bit-bins is evidence in favor of the conjecture.  A
genuinely adversarial integer outside the sample could in principle
violate it.

\paragraph{Connection to the Diophantine literature.}
The cycle-element bound Corollary~\ref{cor:cycle} is essentially a
restatement of a classical observation of Crandall~\cite[Eq.~(8),
p.~1284]{Crandall1978}, who derived for cycles of the unaccelerated
$3x+1$ map the inequality $m_{\min} \leq (3^P - 2^B)^{-1} \cdot
\sum_j 3^{P-1-j} 2^{B_j}$; substituting our bound on the right-hand
sum recovers Corollary~\ref{cor:cycle}.  Eliahou~\cite{Eliahou1993}
combined this with a Diophantine analysis of the continued-fraction
expansion of $\log_2 3$ to prove that any nontrivial cycle has
period $P \geq 17{,}087{,}915$ (his Theorem~3, building on Steiner's
1-cycle elimination~\cite{Steiner1977}).  Our delta is purely
expositional: the bound emerges directly from Lemma~\ref{lem:cycleeq}
and Lemma~\ref{lem:gs} (irrationality), without any
Gelfond--Schneider machinery, which we record because the same
arithmetic ($2^B \neq 3^P$, controlled by the convergents of
$\log_2 3$) drives Theorem~\ref{thm:cf}.

\paragraph{Connection to Tao's almost-everywhere result.}
Tao~\cite{Tao2022} proves that, with respect to the logarithmic
density on $\Nodd$, almost every $n$ has $\liminf_{T \to \infty}
B_T(n)/T \geq \log_2 3$, hence almost every Collatz orbit attains
arbitrarily small ratios $\textit{above}$ $\log_2 3$ in some
finite-$T$ window.  Our complementary contribution is a
\emph{worst-case} (rather than almost-everywhere) statement
restricted to first-descent windows: at every sampled $n$ of
bit-length $\geq 13$ in our experiments the worst-case
\emph{first}-descent ratio is at least $233/147$, and the empirical
floor moves only through the upper continued-fraction objects of
$\log_2 3$.

\paragraph{Open questions.}
Two natural directions are: (i) an effective worst-case lower bound
on $\rho_{\mathrm{des}}(n)$ in terms of $n$, refining
Conjecture~\ref{conj:65-41} to a quantitative form; and (ii) a
Baker-type lower bound for $|B \log 2 - P \log 3|$ specialized to
$(B, P)$ obeying $B/P \geq 65/41$, which would sharpen
Corollary~\ref{cor:cycle}.

\paragraph{Reproducibility.}
All experimental data underlying the tables in Section~\ref{sec:exp}
are documented in the research notebook
\texttt{readme.md} of the repository
\url{https://github.com/weiqi-kids/collatz-research-report}, with an
HTML rendering at \texttt{index.html}; the descent algorithm
(Table~\ref{tab:alg-descent}) is short enough to be re-implemented
in any language with native big-integer arithmetic.  The reported
running times
were measured on a single workstation; total CPU usage across all
experiments is under one CPU-hour.  Verification of the small-$T$
cases of Theorem~\ref{thm:cf} (for example, the claim that
$B(T)/T \geq 65/41$ for $T \leq 93$) reduces to computing
$B(T) = \lfloor T \log_2 3 \rfloor + 1$ for $93$ values of $T$ and
comparing to $65/41$; this can be done by hand or in any computer
algebra system.

%% =================================================================
\section{Conclusion}\label{sec:conc}
%% =================================================================

We have isolated a precise interaction between the continued-fraction
expansion of $\log_2 3$ and the descent dynamics of the compressed
Collatz map.  The deterministic minimum descent ratio $B(T)/T$ for
$T \leq 93$ is governed by the sixth principal convergent
$h_5 = 65/41$ (Theorem~\ref{thm:cf}); an extensive computational
study supports the strong conjecture that this ratio is the actual
worst case across all Collatz orbits at sampled scales
(Conjecture~\ref{conj:65-41}); and the same arithmetic input
(irrationality of $\log_2 3$) yields a clean cycle-element upper
bound (Corollary~\ref{cor:cycle}).  None of these results resolves
the Collatz conjecture, but together they pinpoint the exact
arithmetic obstruction that any future quantitative theory will need
to address.

%% =================================================================
\begin{thebibliography}{99}
%% =================================================================

\bibitem{Baker1990}
A.~Baker,
\emph{Transcendental Number Theory},
2nd ed., Cambridge University Press, Cambridge, 1990.

\bibitem{BarinaCheckup}
D.~Barina,
\emph{Convergence verification of the Collatz problem},
J.\ Supercomput.\ \textbf{77} (2021), 2681--2688.

\bibitem{Bugeaud2004}
Y.~Bugeaud,
\emph{Approximation by Algebraic Numbers},
Cambridge Tracts in Mathematics, vol.~160, Cambridge University
Press, Cambridge, 2004.

\bibitem{ChangPaperA}
L.~Chang,
\emph{A three-state finite automaton for carry propagation in the
accelerated Collatz map: spectral gap and closed-form shift dynamics},
preprint, 2026.

\bibitem{ChangPaperB}
L.~Chang,
\emph{No finite local potential function is a strict Lyapunov
function for the accelerated Collatz map: an impossibility result},
preprint, 2026.

\bibitem{Crandall1978}
R.\,E.~Crandall,
\emph{On the ``$3x+1$'' problem},
Math.\ Comp.\ \textbf{32} (1978), 1281--1292.

\bibitem{Eliahou1993}
S.~Eliahou,
\emph{The $3x+1$ problem: new lower bounds on nontrivial cycle lengths},
Discrete Math.\ \textbf{118} (1993), 45--56.

\bibitem{Gelfond1934}
A.~Gelfond,
\emph{Sur le septi\`eme probl\`eme de Hilbert},
Bull.\ Acad.\ Sci.\ URSS \textbf{7} (1934), 623--634.

\bibitem{HardyWright}
G.\,H.~Hardy and E.\,M.~Wright,
\emph{An Introduction to the Theory of Numbers},
6th ed., Oxford University Press, Oxford, 2008.

\bibitem{Khinchin1964}
A.\,Y.~Khinchin,
\emph{Continued Fractions},
University of Chicago Press, Chicago, 1964.

\bibitem{Lagarias1985}
J.\,C.~Lagarias,
\emph{The $3x+1$ problem and its generalizations},
Amer.\ Math.\ Monthly \textbf{92} (1985), 3--23.

\bibitem{Lagarias2010}
J.\,C.~Lagarias (ed.),
\emph{The Ultimate Challenge: The $3x+1$ Problem},
American Mathematical Society, Providence, RI, 2010.

\bibitem{Schneider1934}
T.~Schneider,
\emph{Transzendenzuntersuchungen periodischer Funktionen, I:
Transzendenz von Potenzen},
J.\ reine angew.\ Math.\ \textbf{172} (1934), 65--69.

\bibitem{Steiner1977}
R.\,P.~Steiner,
\emph{A theorem on the syracuse problem},
in: Proc.\ 7th Manitoba Conf.\ on Numerical Math.\ and Computing,
Congress.\ Numer.\ \textbf{20}, Utilitas Math., Winnipeg, 1977,
pp.\ 553--559.

\bibitem{Tao2022}
T.~Tao,
\emph{Almost all orbits of the Collatz map attain almost bounded
values},
Forum Math.\ Pi \textbf{10} (2022), e12.

\end{thebibliography}

\end{document}
